The growing demand for wireless communications imposes significant challenges related to electromagnetic interference and efficient frequency spectrum allocation. The L(2,1)-Labeling Problem in graphs models the assignment of channels to avoid interference between nearby transmitters, requiring that adjacent vertices have labels with an absolute difference of at least two units, and vertices at distance 2 have an absolute difference of at least one unit. Thus, an L(2,1)-labeling of a graph $G$ is a function $f$ that assigns non-negative integers to the vertex set $V(G)$ such that $|f(u) - f(v)| \geq 2$ if $d(u, v) = 1$ and $|f(u) - f(v)| \geq 1$ if $d(u, v) = 2$. The span of a labeling $f$ is the largest value assigned by $f$ to a vertex of $G$. The objective of the L(2,1)-Labeling problem is to find a valid labeling $f$ that minimizes this value. However, finding such a labeling is an NP-complete problem, which makes exact approaches infeasible for instances of moderate to large size. In this context, this work proposes the investigation of Genetic Algorithms as an efficient heuristic alternative for solving the problem. This work aims to compare the effectiveness of different variations of Genetic Algorithms applied to the L(2,1)-Labeling Problem. The methodology includes the implementation of the algorithms, application to different classes of graphs, parameter tuning, and comparison with an Integer Linear Programming formulation. The results will be analyzed based on solution quality, execution time, and stability. This approach aims to contribute to the optimization of spectral allocation, with relevant impacts on the efficiency of communication networks.

\keywords{L(2,1)-labeling; graphs; genetic algorithms; optimization; heuristics.}